\subsection{Einleitung}
\setauthor{Zoe Emily Öllinger}

Im Rahmen unserer Diplomarbeit bei Hello Again wurde dem Thema Sicherheit und Datenschutz von Beginn an ein hoher Stellenwert gegeben.
Bereits zu Beginn der Diplomarbeit fand eine Einschulung statt, in der uns die Richtlinien zum sicheren Umgang mit Daten, Systemen und Tools erklärt wurden.
Diese Richtlinien betrafen sowohl den Umgang mit sensiblen Kundendaten als auch die erlaubten Tools.

\subsection{Datenschutz und der Umgang mit Kundendaten}
\setauthor{Zoe Emily Öllinger}

Einer der wichtigsten Punkte betraf den Umgang mit Kundendaten.
Uns wurde explizit erklärt, dass Kundendaten zu keinem Zeitpunkt in externe KI-Systeme eingegeben werden dürfen.
Konkret wurde dabei ChatGPT als Beispiel genannt; ein öffentliches KI-Tool, das von OpenAI betrieben wird.
Kundendaten unterliegen in der Europäischen Union der Datenschutz-Grundverordnung (DSGVO), die strenge Anforderungen an die Verarbeitung, Speicherung und Weitergabe personenbezogener Daten stellt.
Werden diese Daten an Dienste weitergegeben, ohne dass eine entsprechende rechtliche Grundlage vorliegt, ist dies ein Datenschutzverstoß, der für das Unternehmen erhebliche rechtliche und finanzielle Konsequenzen haben kann.
Das Vertrauen der Kunden in die Plattform von Hello Again sehr wichtig; jede Form von Databreach würde das Vertrauen beschädigen zwischen Kunde und Firma.
Für uns heißt das, dass wir darauf achten mussten, bei der Nutzung von KI-Tools ausschließlich selbst erstellte Beispieldaten zu verwenden.
Keine echten Kundendaten!


\subsection{Erlaubte KI-Tools: Firmeninterne Lösungen}
\setauthor{Zoe Emily Öllinger}

Hello Again setzt auf firmenintern verwaltete KI-Lösungen.
Im Rahmen der Diplomarbeit durften waren zwei Tools zur Nutzung freigegeben.
Google Gemini (bezahlte Firmenversion): Hello Again verwendet die kostenpflichtige, unternehmensseitig verwaltete Version von Google Gemini.
Der Unterschied zur kostenlosen öffentlichen Version ist, dass Google mit Enterprise-Verträgen garantiert, dass die Daten nicht für das Training von Modellen zu verwendet werden.
Claude als Konsolentool (bezahlte Version im Firmen-Environment): Zusätzlich duften wir Claude von Anthropic verwenden, jedoch ausschließlich in der bezahlten Version.
Claude wurde dabei als Konsolentool eingesetzt, also für Debugging-Unterstützung.

Hier kann man erkennen, dass die Unternehmen die neuartigen KI-Tools nicht ablehnen sondern ihre Mitarbeiter ermutigen diese sinnvoll einzusetzen.


\subsection{Sensibilisierung für Phishing-Angriffe}

Ein weiterer Bestandteil der Einschulung war die Aufklärung über Phishing-Angriffe.
Phishing bezeichnet den Versuch, Mitarbeiter durch gefälschte E-Mails, Links oder Webseiten dazu zu bringen, sensible Informationen herauszugeben zum Beispiel Passwörter, Zugangsdaten uvm.
Für Unternehmen sind Phishing-Angriffe besonders gefährlich, weil sie häufig gezielt auf bestimmte Organisationen abgezielt sind.
Wir wurden gebeten, bei jeder E-Mail mit unbekanntem Absender oder ungewöhnlichem Inhalt besonders Vorsicht zu sein, keine Links einfach so zu öffnen und Rücksprache mit unserem Betreuer zu halten.

\subsection{Branches}
Wir arbeiteten nicht direkt auf dem Main Branch des Projekts.
Der Main Branch enthält den Code der das Dashboard am laufen hält.
Direkte Änderungen an diesem Branch würden ein Risiko darstellen, da fehlerhafte Commits Auswirkungen auf das laufende Dashboard haben könnten bzw haben.
Stattdessen wurde uns ein Branch namens health-monitoring bereitgestellt.
Dieser Branch war unser main Arbeitsbereich.
Er war vom Main Branch isoliert, sodass unsere Änderungen restliche Projekt nicht beeinflussen konnten.
Innerhalb unseres health-monitoring-Branches legten wir uns zwei weitere Branches an, um eine klare Trennung zwischen Frontend und Backend zu gewährleisten, einen frontend-Branch für alle Arbeiten an UI und UX, sowie einen backend-Branch für die serverseitige Logik, API-Anbindungen und Datenbankoperationen.
So konnten effizient Merge-Konflikte reduziert werden.
Am Ende der Diplomarbeit wurde unser erarbeiteter Code nach der Überprüfung, durch die Devs in das offizielle Projekt gemerged.

\subsection{Kein separates Mockup-Projekt}

In vielen Firmen ist es üblich, Praktikanten oder Schülerinnen und Schülern ein Demoprojekt zur Verfügung zu stellen, an dem sie arbeiten können.
In unserem Fall war dies firmenintern nicht machbar.
Der Grund war personeller und zeitlicher Aufwand.
Das Aufsetzen eines Mockup-Projekts hätte viel Zeit und Entwickler gekostet.
Da diese Devs nicht verfügbar waren, entschied man sich für einen eigenen Branch.
Diese Entscheidung hatte sowohl Vorteile als auch Nachteile.
Der offensichtliche Vorteil liegt darin, dass wir mit echten firmeninternen Strukturen gearbeitet haben aber der Nachteil besteht darin, dass die Fehlertoleranz geringer ist.
Schwerwiegende Fehler im Branch hätten den Merge-Prozess erschwert.
Da aber die Entwickler unsere Commits regelmäßig begutachtet haben, hatten wir diese Probleme nicht.








































